\documentclass[11pt]{article}
\usepackage{amsmath,amssymb,amsthm,geometry}
\geometry{margin=1in}
\title{A counterexample to Conjecture 00000002304}
\author{gaochengzhecpu (AI-assisted submission)}
\date{October 3, 2026}
\setlength{\emergencystretch}{3em}
\begin{document}
\maketitle

The source proposes an equality classification under \emph{general
linear actions}: equality in $k(GV)\leq k(G)k(V)$ occurs exactly when,
among other conditions, the $G$-module $V$ is faithful.
Here $k(H)$ denotes the number of conjugacy classes of a finite group
$H$, and $k(V)$ refers to the additive group of the vector space.
We give a general linear action attaining equality that is not faithful.
Both language versions state the general-action scope; neither imposes
faithfulness or coprimality as a hypothesis.

\paragraph{Construction.}
Let $G=C_2=\{0,1\}$ with addition modulo $2$ and identity $0$.
Let the field be $\mathbb F_2=\{0,1\}$, with its usual addition and
multiplication modulo $2$, and put $V=\mathbb F_2$, viewed as a
one-dimensional vector space over itself.
Define the action by
\[
 \rho(g)(v)=v\qquad(g\in G,\ v\in V).
\]
Every action map is the identity linear automorphism; the identity and
composition laws for a group action hold. Thus this is a legitimate
linear action. Its kernel is all of the nontrivial group $G$, so it is
not faithful. Equivalently, $\rho(0)=\rho(1)$ while $0\ne1$.
The vector space is nonzero and one-dimensional; in particular, the
failure of faithfulness is not an artifact of using the zero module.

Use the standard semidirect product on the carrier $V\times G$:
\[
 (v,g)(w,h)=(v+\rho(g)(w),\,g+h)=(v+w,\,g+h).
\]
Consequently $V\rtimes_\rho G=C_2\times C_2$, the four-element
abelian group. Its identity is $(0,0)$ and each element is its own
inverse.

\paragraph{Conjugacy classes and contradiction.}
In an abelian group, conjugating $x$ by any $g$ gives
$gxg^{-1}=x$. Hence each element is a singleton conjugacy class.
Therefore
\[
 k(G)=2,\qquad k(V)=2,\qquad k(V\rtimes_\rho G)=4=2\cdot2=k(G)k(V).
\]
Equality holds, yet $V$ is not a faithful $G$-module. This contradicts
the faithfulness condition required by the proposed classification.
It also contradicts the claim of strict inequality for all remaining
actions. No interpretation of the additional character-degree
assertion can restore this false necessary condition.

\paragraph{Scope.}
This disproof addresses the conjecture exactly as written, which
quantifies over general linear actions. One could ask a different
question after imposing a faithful-action hypothesis, or further
hypotheses such as coprimality. Our example is deliberately not faithful,
and $|G|=2$ equals the characteristic of the field, so it is not offered
as a counterexample to either of those restricted variants.
The source's stated equality classification cannot itself supply a
missing assumption excluding counterexamples to that classification.

\paragraph{Formal proof and correspondence with the mathematical objects.}
The file \texttt{lean/Main.lean} is self-contained over \texttt{Std},
using Lean 4.19.0.

\texttt{GroupData} contains multiplication, identity, inverse, and
all group laws. \texttt{FieldData} contains the additive and
multiplicative field laws, distributivity, inverses of all nonzero
elements, and the axiom $0\ne1$. \texttt{VectorSpaceData} contains an
abelian additive group and the vector-space scalar laws.
The concrete values \texttt{C2}, \texttt{F2}, and \texttt{V2}
verify these axioms with the modulo-$2$ operations on \texttt{Fin 2}.
The results \texttt{every\_vector\_is\_scalar} and
\texttt{basis\_vector\_nonzero} additionally exhibit the nonzero
spanning vector $1$.

\texttt{LinearAction} asserts the group action laws and preservation
of addition and scalar multiplication.
The theorem \texttt{action\_inverse} records the inverse-action identity,
so the acting maps are indeed automorphisms.
\texttt{trivialAction} is the action used above.
\texttt{Faithful} is injectivity of the map from group elements to
action maps, and \texttt{not\_faithful} uses the distinct elements
$0$ and $1$ directly.

The multiplication of the group \texttt{GV}, defined on the entire
product \texttt{Fin 2}$\times$\texttt{Fin 2}, is precisely the displayed
semidirect formula. All group laws are verified; the theorems
\texttt{semidirect\_multiplication} and
\texttt{semidirect\_identity} expose that correspondence.

In particular, the formalization does not replace conjugacy-class
counts by unrelated constants.
\texttt{Conjugate H x y} is defined by the existence of a group element
$g$ with $gxg^{-1}=y$.
\texttt{HasClassNumber H n} requires a setoid with exactly this
relation and a bijection from \texttt{Fin n} onto its actual
\texttt{Quotient}.
The formal proof verifies that conjugacy is equality in each of the
three groups. It then builds explicit quotient bijections, including
a complete four-element enumeration of the product carrier, proving
the class numbers $2,2,4$.

\texttt{ClaimedEqualityNecessity} is the faithfulness implication
required by the source, restricted to the specific finite group, field,
and vector space already constructed, but quantified over their linear
actions and corresponding semidirect products.
A universal general-action equality classification necessarily implies
this restricted claim.
\texttt{counterexample} states the class-number equalities and failure
of faithfulness together; \texttt{conjecture2304\_false} proves the
negation of that necessary implication.

The proof uses kernel-reduced \texttt{decide}, with no
\texttt{native\_decide}, admitted statement, or custom axiom.
Printed axioms of the final theorem are exactly
\texttt{propext}, \texttt{Classical.choice}, and \texttt{Quot.sound}.
\end{document}
